%----------------------------------------------------------------------------------------
% NOTATIE EN FIGUREN VAN DE CURSUS NUMERIEKE METHODES
% Staat in de preambule en niet in een hoofdstuk: ook de lijsten met oplossingen gebruiken dit.
%----------------------------------------------------------------------------------------
% Bewaarde grafiek, geplaatst door maak-pdf.py: bestand en pagina in figuren-cache/
\newcommand{\grafiekbestand}[2]{\includegraphics[page=#2,trim=3bp 3bp 3bp 3bp]{figuren-cache/#1}}

\providecommand{\N}{\mathbb{N}}
\providecommand{\R}{\mathbb{R}}
\newcommand{\dx}{\,\mathrm{d}x}
\newcommand{\dt}{\,\mathrm{d}t}
\newcommand{\e}{\mathrm{e}}
\newcommand{\ig}{\displaystyle\int}
\newcommand{\bgtan}{\operatorname{bgtan}}
\newcommand{\grens}[3]{\Bigl[#1\Bigr]_{#2}^{#3}}
% Rechter- en linkerlimiet, met de ongelijkheid onder de pijl
\newcommand{\limr}[2]{\lim\limits_{\substack{#1 \to #2 \\ >}}}
\newcommand{\liml}[2]{\lim\limits_{\substack{#1 \to #2 \\ <}}}
% Benadering met k stappen of deelintervallen: M_n, T_n, S_n
\newcommand{\Mn}[1]{M_{#1}}
\newcommand{\Tn}[1]{T_{#1}}
\newcommand{\Sn}[1]{S_{#1}}
% Een iteratie: x_{n+1} = g(x_n). Bovenindex tussen haakjes voor de stap bij stelsels: x^{(k)}
\newcommand{\stap}[2]{#1^{(#2)}}

% Zoals stappen, maar met een bredere kolom voor de berekening: voor lange formules
\newenvironment{stappenbreed}{%
  \par\renewcommand{\arraystretch}{1.25}%
  \noindent\begin{tabular}{@{}>{$\displaystyle}p{9.8cm}<{$}@{\hspace{4mm}}>{\RaggedRight\small\itshape\arraybackslash}p{4.9cm}@{}}%
}{\end{tabular}\par}

% Zoals stappen, met eigen kolombreedtes: \begin{stappenvrij}{6.2cm}{3.8cm}. Voor naast een figuur.
\newenvironment{stappenvrij}[2]{%
  \par\renewcommand{\arraystretch}{1.25}%
  \noindent\begin{tabular}{@{}>{$\displaystyle}p{#1}<{$}@{\hspace{3mm}}>{\RaggedRight\small\itshape\arraybackslash}p{#2}@{}}%
}{\end{tabular}\par}

% Tekst of berekening links, figuur rechts: \naast[breedte rechts]{links}{rechts}
% De breedte rechts komt uit de figuur zelf; het optionele argument wordt niet gebruikt.
\newsavebox{\naastbox}
\newcommand{\naast}[3][]{%
  \par\noindent\sbox{\naastbox}{#3}%
  \begin{minipage}[t]{\dimexpr\linewidth-\wd\naastbox-3mm\relax}\vspace{0pt}#2\end{minipage}\hfill
  \begin{minipage}[t]{\wd\naastbox}\vspace{0pt}\usebox{\naastbox}\end{minipage}\par\addvspace{1mm}}
% Een rij figuren, gecentreerd
\newcommand{\figuren}[1]{\par\addvspace{1mm}{\centering #1\par}\addvspace{1.5mm}}
% Opgave met een grafiek rechts: \begin{metgrafiek}[breedte]{tikz} opgave en deelvragen \end{metgrafiek}
\newsavebox{\grafiekbox}
\newenvironment{metgrafiek}[2][]{%
  \sbox{\grafiekbox}{#2}%
  \begin{minipage}[t]{\dimexpr\linewidth-\wd\grafiekbox-3mm\relax}%
}{%
  \end{minipage}\hfill
  \begin{minipage}[t]{\wd\grafiekbox}\vspace{-2mm}\usebox{\grafiekbox}\end{minipage}\par}

% Tabel met een raster, alles gecentreerd: \begin{rastertabel}{ccc} ... \end{rastertabel}
\newenvironment{rastertabel}[1]{\par\addvspace{1mm}\begingroup\centering\renewcommand{\arraystretch}{1.3}%
  \setlength{\arrayrulewidth}{.4pt}\arrayrulecolor{lijngrijs}\begin{tabular}{#1}}%
  {\end{tabular}\par\endgroup\addvspace{1.5mm}}

% Rooster met assen: \rooster{xmin}{xstap}{xmax}{ymin}{ystap}{ymax}, met getallen op de assen (gehele waarden).
% \roosterleeg doet hetzelfde zonder getallen, voor eigen labels.
\newcommand{\roosterleeg}[6]{%
  \draw[black!14,very thin,xstep=#2,ystep=#5] (#1,#4) grid (#3,#6);
  \draw[->,grijs] (#1,0) -- ({#3+.4*#2},0) node[right,text=black] {$x$};
  \draw[->,grijs] (0,#4) -- (0,{#6+.4*#5}) node[above,text=black] {$y$};}
\newcommand{\rooster}[6]{%
  \roosterleeg{#1}{#2}{#3}{#4}{#5}{#6}
  \pgfmathtruncatemacro{\rtweex}{#1+#2}\pgfmathtruncatemacro{\rtweey}{#4+#5}
  \foreach \t in {#1,\rtweex,...,#3}{\ifnum\t=0\else\node[below,inner sep=1.5pt,text=black] at (\t,0) {$\t$};\fi}
  \foreach \t in {#4,\rtweey,...,#6}{\ifnum\t=0\else\node[left,inner sep=1.5pt,text=black] at (0,\t) {$\t$};\fi}}
% Assen zonder rooster
\newcommand{\assen}[4]{%
  \draw[->,grijs] (#1,0) -- (#2,0) node[right,text=black] {$x$};
  \draw[->,grijs] (0,#3) -- (0,#4) node[above,text=black] {$y$};}
% In de functievoorschriften is \x de veranderlijke; schrijf (\x)^2 en sin(\x r).
\newcommand{\kromme}[4][]{\draw[accent,thick,#1] plot[domain=#2:#3,samples=90,smooth] (\x,{#4});}
% \gebied: tussen de grafiek en de horizontale as
\newcommand{\gebied}[4][accent!30]{\fill[#1] plot[domain=#2:#3,samples=70] (\x,{#4}) -- (#3,0) -- (#2,0) -- cycle;}
% \riemann{aantal}{a}{b}{voorschrift}{plaats}: rechthoeken, met de hoogte op plaats 0 (links), 0.5 (midden) of 1 (rechts)
\newcommand{\riemann}[5]{%
  \pgfmathsetmacro{\rbreedte}{(#3-(#2))/#1}%
  \foreach \i in {1,...,#1}{%
    \pgfmathsetmacro{\xl}{#2+(\i-1)*\rbreedte}%
    \pgfmathsetmacro{\xr}{\xl+\rbreedte}%
    \pgfmathsetmacro{\x}{\xl+#5*\rbreedte}%
    \pgfmathsetmacro{\hoogte}{#4}%
    \draw[accent,fill=accent!22] (\xl,0) rectangle (\xr,\hoogte);}}
% \trapezia{aantal}{a}{b}{voorschrift}: de trapezia van de trapeziumregel
\newcommand{\trapezia}[4]{%
  \pgfmathsetmacro{\rbreedte}{(#3-(#2))/#1}%
  \foreach \i in {1,...,#1}{%
    \pgfmathsetmacro{\xl}{#2+(\i-1)*\rbreedte}%
    \pgfmathsetmacro{\xr}{\xl+\rbreedte}%
    \pgfmathsetmacro{\x}{\xl}\pgfmathsetmacro{\yl}{#4}%
    \pgfmathsetmacro{\x}{\xr}\pgfmathsetmacro{\yr}{#4}%
    \draw[accent,fill=accent!22] (\xl,0) -- (\xl,\yl) -- (\xr,\yr) -- (\xr,0) -- cycle;}}

% Webgrafiek van x_{n+1} = g(x_n): \webgrafiek[opties]{g met \x}{x0}{aantal stappen}
% Begint op de x-as in x0, gaat verticaal naar de grafiek van g en horizontaal naar y = x.
\newcommand{\webgrafiek}[4][]{%
  \xdef\webx{#3}%
  \foreach \k in {1,...,#4}{%
    \pgfmathsetmacro{\x}{\webx}\pgfmathsetmacro{\weby}{#2}%
    \ifnum\k=1 \draw[fel!80!black,thick,#1] (\webx,0) -- (\webx,\weby);
    \else \draw[fel!80!black,thick,#1] (\webx,\webx) -- (\webx,\weby);\fi
    \draw[fel!80!black,thick,#1] (\webx,\weby) -- (\weby,\weby);
    \xdef\webx{\weby}}}
% Methode van Newton: \newtonstappen{f met \x}{f' met \x}{x0}{aantal stappen}
% Per stap: verticaal van (x_n, 0) naar de grafiek, en langs de raaklijn naar (x_{n+1}, 0).
\newcommand{\newtonstappen}[4]{%
  \xdef\nwx{#3}%
  \foreach \k in {1,...,#4}{%
    \pgfmathsetmacro{\x}{\nwx}\pgfmathsetmacro{\nwf}{#1}\pgfmathsetmacro{\nwa}{#2}%
    \pgfmathsetmacro{\nwv}{\nwx-\nwf/\nwa}%
    \draw[grijs,dashed] (\nwx,0) -- (\nwx,\nwf);
    \draw[fel!80!black,thick] (\nwx,\nwf) -- (\nwv,0);
    \fill[accent] (\nwx,\nwf) circle (1.5pt);
    \fill[black] (\nwv,0) circle (1.3pt);
    \xdef\nwx{\nwv}}}

% Programma's (thuis, met de computer): \begin{python} ... \end{python}, niet in een voorbeeld of oefening
\usepackage{listings}
\lstdefinestyle{py}{language=Python,basicstyle=\ttfamily\small,keywordstyle=\color{accent}\bfseries,
  commentstyle=\color{grijs}\itshape,stringstyle=\color{inkt},columns=fullflexible,keepspaces=true,
  showstringspaces=false,upquote=true,aboveskip=0pt,belowskip=0pt}
\lstnewenvironment{python}{\lstset{style=py,backgroundcolor=\color{black!4},frame=leftline,framerule=1pt,
  rulecolor=\color{grijs},xleftmargin=4mm,framexleftmargin=3mm,framextopmargin=1mm,framexbottommargin=1mm,
  aboveskip=1.5ex,belowskip=1.5ex}}{}
